\documentclass[10.5pt, a4paper, twoside, openright]{report}

% -- English --
%\usepackage[english]{babel}
% -- French  --
\usepackage[francais]{babel}

\usepackage[utf8x]{inputenc}
\usepackage[a4paper, left=2cm, right=2cm, top=2.5cm, bottom=2cm, headheight=40pt]{geometry}
\usepackage{graphicx}

% -- Layout --
\usepackage{emptypage}
\usepackage{fancyhdr}
\pagestyle{fancy}
\renewcommand{\headrulewidth}{1pt}
\renewcommand{\chaptermark}[1]{\markboth{#1}{}}
\fancyhead[C]{\leftmark}
\fancyhead[RE,LO]{}
\fancyhead[RO,LE]{\includegraphics[scale=0.016]{UB.jpg}}
\newcommand{\HRule}{\rule{\linewidth}{0.5mm}}
\raggedbottom{}

\title{Analyse de malware statique et dynamique : Les outils et quelques cas pratiques}
\author{Erwan Grelet \& Amélie Guémon}

\begin{document}

% -- Title Page --
\begin{titlepage}
  \begin{sffamily}
  \begin{center}

    Univesity of Bordeaux\\College of Science \& Technology\\
           351 Cours de la Liberation\\33400 Talence\\[1em]
            \textbf{\underline{- Projet M2 -}}\\[1.5cm]

    \includegraphics[scale=0.11]{UB.jpg}
    \\[3cm]

    \HRule\\[0.3cm]
    { \huge \bfseries Analyse de malware statique et dynamique : \\Les outils et quelques cas pratiques\\[0.5cm] }
    \HRule\\
    \begin{flushright}
      \bfseries {Erwan GRELET \& Amélie GUÉMON\\[1em]Master 2 CSI --- Cryptology \& Computer Security}\\[6cm]
    \end{flushright}

    \today

  \end{center}
  \end{sffamily}
\end{titlepage}

\tableofcontents

Coucou à toi !

% -- New Part --
%\include{tex\part}

% -- For Biblio --
%\nocite{*}
%\bibliographystyle{plain}
%\bibliography{biblio}

\end{document}
